% ============================================================
% SECCIÓN 9 — Conclusiones
% ============================================================

\section{Conclusiones}

% --- Diapositiva 1: La cadena de resultados ---
\begin{frame}{El Arco de la Teoría de Ramsey}

    \begin{center}
    \begin{tikzpicture}[
        node distance=0.6cm and 0.4cm,
        boxa/.style={draw, rounded corners=4pt, fill=blue!60,   text=white,
                     font=\small\bfseries, minimum width=2.4cm,
                     minimum height=0.7cm, align=center},
        boxb/.style={draw, rounded corners=4pt, fill=teal!65,   text=white,
                     font=\small\bfseries, minimum width=2.4cm,
                     minimum height=0.7cm, align=center},
        boxc/.style={draw, rounded corners=4pt, fill=violet!65, text=white,
                     font=\small\bfseries, minimum width=2.4cm,
                     minimum height=0.7cm, align=center},
        boxd/.style={draw, rounded corners=4pt, fill=orange!70, text=white,
                     font=\small\bfseries, minimum width=2.4cm,
                     minimum height=0.7cm, align=center},
        boxe/.style={draw, rounded corners=4pt, fill=red!65,    text=white,
                     font=\small\bfseries, minimum width=2.4cm,
                     minimum height=0.7cm, align=center},
        arr/.style={->, thick, gray!70}
    ]
        \node[boxa] (ramsey) {Teoría de\\Ramsey};
        \node[boxb, right=0.5cm of ramsey] (schur)  {Schur\\(1916)};
        \node[boxc, right=0.5cm of schur]  (vdw)    {Van der\\Waerden (1927)};
        \node[boxd, below=0.8cm of schur]  (szem)   {Szemerédi\\(1975)};
        \node[boxe, right=0.5cm of szem]   (gt)     {Green-Tao\\(2004)};

        \draw[arr] (ramsey) -- (schur);
        \draw[arr] (schur)  -- (vdw);
        \draw[arr] (vdw)    -- (szem);
        \draw[arr] (szem)   -- (gt);

        \node[font=\tiny, gray, below=0.15cm of ramsey] {patrones inevitables};
        \node[font=\tiny, gray, below=0.15cm of schur]  {$x{+}y{=}z$ mono.};
        \node[font=\tiny, gray, below=0.15cm of vdw]    {PA mono.};
        \node[font=\tiny, gray, below=0.15cm of szem]   {$d{>}0 \to$ PA};
        \node[font=\tiny, gray, below=0.15cm of gt]     {primos $\to$ PA};
    \end{tikzpicture}
    \end{center}

    \vspace{0.2cm}

    \begin{block}{El mensaje central}
        Cada resultado extiende el anterior, forzando más orden
        en contextos donde parecía no haberlo.
        \textbf{La estructura aritmética es inevitable.}
    \end{block}

\end{frame}


% --- Diapositiva 2: Reflexión final ---
\begin{frame}{Conclusiones}

    \begin{exampleblock}{\centering Principio de Motzkin — revisitado}
        \centering
        \Large\textit{``El desorden completo es imposible.''}
    \end{exampleblock}

    \vspace{0.35cm}

    \begin{columns}[T]
        \begin{column}{0.48\textwidth}
            \textbf{Lo que aprendimos:}
            \begin{itemize}
                \item El \textbf{Palomar} es el germen de todo.
                \item \textbf{Schur} y \textbf{VdW} garantizan
                      estructuras en coloraciones.
                \item \textbf{Szemerédi} solo necesita densidad positiva.
                \item \textbf{Green-Tao} vence incluso a la
                      densidad nula de los primos.
            \end{itemize}
        \end{column}
        \begin{column}{0.48\textwidth}
            \textbf{Preguntas abiertas:}
            \begin{itemize}
                \item ¿Cuánto vale $R(5,5)$?
                \item Cotas exactas de $W(r,k)$.
                \item ¿PA en los primos gemelos?
                \item Conjetura de Erdős sobre PA en
                      conjuntos de series divergentes.
            \end{itemize}
        \end{column}
    \end{columns}

    \pause
    \vspace{0.35cm}

    \begin{alertblock}{\centering Mensaje final}
        \centering
        En matemáticas, los \textbf{patrones no son casualidad}:
        son una consecuencia \textbf{inevitable} del tamaño.
    \end{alertblock}

\end{frame}
